\documentclass[11pt]{article}
\usepackage[T1]{fontenc}
\usepackage{textcomp}
\usepackage[margin=0.75in]{geometry}
\usepackage{amsmath,amssymb,amsthm}
\usepackage{array,booktabs}
\usepackage{hyperref}
\setlength{\parindent}{0pt}
\newtheorem{theorem}{Theorem}
\theoremstyle{definition}
\newtheorem{definition}{Definition}
\title{Flow Proof Artifact: prop-09-inscribe-pentagon}
\author{Generated by \texttt{flow doc proof}}
\date{Flow Proof Book}
\begin{document}
\maketitle
In a given circle to inscribe a regular pentagon.\\[0.5em]
\textbf{Source.} euclid — Elements, Book IV, Proposition 9\\[0.5em]
\bigskip
\noindent{\large \textbf{Derived fact 1.} \textit{In a given circle to inscribe a regular pentagon} \hfill the Euclidean plane \cdot Euclid Book IV \cdot Proposition 9: in a circle to inscribe a regular pentagon}\par
\smallskip\noindent\textit{Coordinate: the Euclidean plane \cdot Euclid Book IV \cdot Proposition 9: in a circle to inscribe a regular pentagon}\par
\smallskip\noindent\textit{Source: euclid — Elements, Book IV, Proposition 9}\par
\smallskip\noindent\textit{Built on: proposition 10: to construct an isosceles triangle with base angles double the vertex, for Euclid Book IV on the Euclidean plane, proposition 20: the central angle is double the inscribed angle on the same arc, for Euclid Book III on the Euclidean plane}\par
\medskip
\noindent\textbf{Goal.} In a given circle to inscribe a regular pentagon\par
\noindent$\displaystyle \text{regular pentagon inscribed}$\par
\medskip
\noindent
\begin{tabular}{@{} >{\raggedright\arraybackslash}p{0.06\textwidth} >{\raggedright\arraybackslash}p{0.47\textwidth} >{\raggedleft\arraybackslash}p{0.41\textwidth} @{}}
\textbf{\#} & \textbf{Proof} & \textbf{Mathematics} \\
\midrule
\textbf{1.} & We prove that proposition 9: in a circle to inscribe a regular pentagon for Euclid Book IV the Euclidean plane. & \\[0.45em]
\textbf{2.} & Let circle = given\_circle. & \\[0.45em]
\textbf{3.} & Let pentagon = inscribed\_regular\_pentagon. & \\[0.45em]
\textbf{4.} & Let side = pentagon\_side. & \\[0.45em]
\textbf{5.} & From step 2, step 3, and step 4, we can deduce that regular pentagon in circle. & $\displaystyle \text{regular pentagon in circle}$ \\[0.45em]
\textbf{6.} & We invoke the derived fact governing Euclid Book IV the Euclidean plane: proposition 10: to construct an isosceles triangle with base angles double the vertex, for Euclid Book IV on the Euclidean plane. & \\[0.45em]
\textbf{7.} & We invoke the derived fact governing Euclid Book III the Euclidean plane: proposition 20: the central angle is double the inscribed angle on the same arc, for Euclid Book III on the Euclidean plane. & \\[0.45em]
\textbf{8.} & From step 6 and step 7, this implies regular pentagon inscribed. Hence proven. & $\displaystyle \text{regular pentagon inscribed}$ \\[0.45em]
\bottomrule
\end{tabular}
\label{thm:Geometry-Euclid Book IV-proposition_9_in_a_circle_to_inscribe_a_regular_pentagon}
\par\medskip\hrule
\end{document}
